\chapter{التناسب والسرعة والمتوسطات}\label{ch:g8:speed}

يلتقي \omterm{def:g7:prop:table}{التناسب} (\cref{ch:g7:prop}) الآن العالمَ الفيزيائي:
\omterm{def:g8:speed:def}{السرعة} والصبيب وتحويلات الوحدات --- وهي مقادير
\emph{خوارج قسمة} مقدارين آخرين. وينتهي الفصل بالمتوسطات
المرجّحة، حيث يمنع التفكير التناسبي خطأً كلاسيكيًّا.

\section{السرعة المتوسطة}

\begin{definition}[السرعة المتوسطة]\label{def:g8:speed:def}
\emph{السرعة المتوسطة}\index{سرعة} لرحلة هي \omterm{def:g3:division:remainder}{خارج القسمة}
\[
v = \frac{d}{t}
\qquad
\left(\text{المسافة مقسومة على المدّة}\right),
\]
بالكيلومتر في الساعة (كم/سا)، أو المتر في الثانية (م/ث)، \dots\ و
الصيغة تنقلب: $d = v \times t$ و $t = \frac{d}{v}$.
\end{definition}

\begin{example}\label{ex:g8:speed:basic}
يقطع قطار $270$ كم في $1$ سا و $30$ دق. وأولًا حوّل الزمن:
$1$ سا و $30$ دق $= 1.5$ سا. ثم
\[
v = \frac{270}{1.5} = 180 \text{ km/h}.
\]
والمسافة \omterm{def:g8:speed:def}{بسرعة} $90$ كم/سا خلال $2$ سا و $20$ دق $= \frac73$ سا:
$d = 90 \times \frac73 = 210$ كم. والزمن اللازم لقطع $35$ كم \omterm{def:g8:speed:def}{بسرعة} $14$ كم/سا:
$t = \frac{35}{14} = 2.5$ سا $= 2$ سا و $30$ دق.
\end{example}

\begin{remark}[الدقائق ليست أعدادًا عشرية]
المدّة $1$ سا و $30$ دق هي $1.5$ سا، لكن $1$ سا و $20$ دق \emph{ليست} $1.2$ سا:
بل هي $1 + \frac{20}{60} = \frac43 \approx 1.33$ سا. فحوّل دائمًا
الدقائق إلى \omterm{def:g6:fractions:def}{كسر} من الساعة ($60$ دق $= 1$ سا) قبل القسمة.
\end{remark}

\begin{omfigure}
\begin{tikzpicture}
\begin{axis}[omaxis,
  width=8.4cm, height=5.6cm,
  xmin=0, xmax=3.4, ymin=0, ymax=270,
  xlabel={الزمن (سا)}, ylabel={المسافة (كم)},
  xtick={1,2,3}, ytick={80,160,240}]
  \addplot[omDef, thick, domain=0:3.2] {80*x};
  \addplot[omThm, thick, domain=0:3.2] {50*x};
  \node[omDef, font=\small, rotate=36] at (axis cs:2.35,215)
    {$80$ كم/سا};
  \node[omThm, font=\small, rotate=25] at (axis cs:2.6,155)
    {$50$ كم/سا};
\end{axis}
\end{tikzpicture}

{\small \omterm{def:g8:speed:def}{بسرعة} ثابتة، تكون المسافة متناسبة مع الزمن: أي
\omterm{def:g6:lines:objects}{مستقيم} يمرّ بالمبدأ، وميله \emph{هو}
\omterm{def:g8:speed:def}{السرعة}.}
\end{omfigure}

\begin{method}[تحويل السرعات]\label{met:g8:speed:convert}
لتحويل كم/سا إلى م/ث: $1$ كم $= 1000$ م و $1$ سا $= 3600$ ث، إذن
\[
v \text{ كم/سا} = \frac{v \times 1000}{3600} \text{ م/ث}
= \frac{v}{3.6} \text{ م/ث}.
\]
ولتحويل م/ث إلى كم/سا، اِضرب في $3.6$.
\end{method}

\begin{example}\label{ex:g8:speed:convert}
$90$ كم/سا $= \frac{90}{3.6} = 25$ م/ث. والعدّاء الذي يجري $100$ م في
$10$ ث يتحرّك \omterm{def:g8:speed:def}{بسرعة} $10$ م/ث $= 36$ كم/سا.
\end{example}

\section{مقادير خوارج القسمة}

\begin{example}[الصبيب والكتلة الحجمية وثمن الكيلوغرام]\label{ex:g8:speed:quotients}
\omterm{def:g8:speed:def}{للسرعة} أبناء عمّ كثيرون، تُعالَج كلها بالطريقة نفسها:
\begin{itemize}
  \item حنفية تملأ $48$ لترًا في $4$ دق: فالصبيب
        $\frac{48}{4} = 12$ لتر/دق، ومنه فملء حوض سعته $150$ لترًا يستغرق
        $\frac{150}{12} = 12.5$ دق؛
  \item و $0.6$ كغ من الجبن ثمنها $9$ يورو: فالثمن
        $\frac{9}{0.6} = 15$ يورو للكيلوغرام؛
  \item وسيارة تستهلك $6.3$ لترًا لكل $90$ كم: فالاستهلاك
        $\frac{6.3}{90} \times 100 = 7$ لترات لكل $100$ كم.
\end{itemize}
حدّد \omterm{def:g3:division:remainder}{خارج القسمة}، وتقوم الصيغة الثلاثية \omterm{def:g3:division:remainder}{بالباقي}
($q = \frac ab$، و $a = qb$، و $b = \frac aq$).
\end{example}

\section{المتوسطات المرجّحة}

\begin{definition}[المتوسط المرجّح]\label{def:g8:speed:average}
حين تأتي القيم $v_1, v_2, \dots$ مع تعدادات (أو أوزان)
$n_1, n_2, \dots$، يكون \emph{متوسطها المرجّح}\index{متوسط مرجح}
هو
\[
\bar v = \frac{n_1 v_1 + n_2 v_2 + \dots}{n_1 + n_2 + \dots} :
\]
أي \omterm{def:g1:addition:def}{مجموع} القيم، مقسومًا على \omterm{def:g1:addition:def}{مجموع} الأوزان.
\end{definition}

\begin{example}\label{ex:g8:speed:average}
اِمتحان أجرته مجموعتان: المجموعة 1 فيها $12$ تلميذًا ومتوسطها
$11$؛ والمجموعة 2 فيها $18$ تلميذًا ومتوسطها $14$. ومتوسط القسم
كله:
\[
\bar v = \frac{12 \times 11 + 18 \times 14}{12 + 18}
= \frac{132 + 252}{30} = \frac{384}{30} = 12.8 .
\]
و\emph{ليس} $\frac{11 + 14}{2} = 12.5$: فالمجموعة الأكبر تجذب
المتوسط نحو متوسطها --- والأوزان مهمّة.
\end{example}

\begin{example}[السرعة المتوسطة على مرحلتين]\label{ex:g8:speed:twolegs}
يقطع دراج $30$ كم \omterm{def:g8:speed:def}{بسرعة} $30$ كم/سا، ثم $30$ كم \omterm{def:g8:speed:def}{بسرعة} $15$ كم/سا. و
\omterm{def:g8:speed:def}{السرعة المتوسطة} على الرحلة كلها \emph{ليست}
$\frac{30 + 15}{2} = 22.5$ كم/سا. فاِحسب بالتعريف:
\begin{enumerate}
  \item الزمنان: $\frac{30}{30} = 1$ سا، ثم $\frac{30}{15} = 2$ سا؛
        \omterm{def:g1:addition:def}{والمجموع} $3$ سا؛
  \item والمسافة الكلية $60$ كم، إذن
        $v = \frac{60}{3} = 20$ كم/سا.
\end{enumerate}
فالمرحلة البطيئة تدوم أطول، ومنه فوزنها أكبر --- فمتوسط السرعات
يجب أن يُرجَّح \emph{بالزمن}.
\end{example}

\section{تمارين}

\begin{exercise}[$\star$]\label{exo:g8:speed:1}
اِحسب \omterm{def:g8:speed:def}{السرعة المتوسطة}: $150$ كم في $2$ سا؛ و $27$ كم في $45$ دق؛
و $100$ م في $12.5$ ث.
\end{exercise}

\begin{exercise}[$\star$]\label{exo:g8:speed:2}
اِحسب المسافة: $2$ سا \omterm{def:g8:speed:def}{بسرعة} $85$ كم/سا؛ و $40$ دق \omterm{def:g8:speed:def}{بسرعة} $90$ كم/سا.
واِحسب الزمن: $315$ كم \omterm{def:g8:speed:def}{بسرعة} $90$ كم/سا (بالساعة والدقيقة).
\end{exercise}

\begin{exercise}[$\star$]\label{exo:g8:speed:3}
حوّل: $72$ كم/سا إلى م/ث؛ و $5$ م/ث إلى كم/سا؛ و $108$ كم/سا إلى م/ث.
\end{exercise}

\begin{exercise}[$\star$]\label{exo:g8:speed:4}
حنفية تعطي $15$ لتر/دق. فكم من الوقت لملء خزّان سعته $600$ لتر؟ وكم
لترًا في $2$ سا و $30$؟
\end{exercise}

\begin{exercise}[$\star$]\label{exo:g8:speed:5}
أيّهما صفقة أفضل: $1.2$ كغ من التفاح بمبلغ $3.30$ يورو، أم
$0.8$ كغ بمبلغ $2.16$ يورو؟ قارن الأثمان للكيلوغرام.
\end{exercise}

\begin{exercise}[$\star$]\label{exo:g8:speed:6}
قسم فيه $25$ تلميذًا متوسطه $12.4$ في اِختبار؛ وقسم آخر
فيه $35$ تلميذًا متوسطه $10$. اِحسب متوسط
القسمين معًا.
\end{exercise}

\begin{exercise}[$\star$]\label{exo:g8:speed:7}
ينتقل الصوت \omterm{def:g8:speed:def}{بسرعة} نحو $340$ م/ث. وترى ومضة برق وتسمع
الرعد بعدها بمدّة $6$ ثوانٍ. فكم بَعُد موضع سقوط البرق
(إلى أقرب $100$ م)؟
\end{exercise}

\begin{exercise}[$\star\star$]\label{exo:g8:speed:8}
متجوّل يمشي $2$ سا \omterm{def:g8:speed:def}{بسرعة} $5$ كم/سا، ويستريح $30$ دق، ثم يمشي $1$ سا و $30$
\omterm{def:g8:speed:def}{بسرعة} $4$ كم/سا. اِحسب المسافة الكلية \omterm{def:g8:speed:def}{والسرعة المتوسطة}
\emph{بما فيها الاستراحة} (المسافة الكلية على الزمن الكلي المنقضي).
\end{exercise}

\begin{exercise}[$\star\star$]\label{exo:g8:speed:9}
علامات بمعاملات: حصلت تلميذة على $15$ (بمعامل $3$)، و $9$
(بمعامل $2$)، و $12$ (بمعامل $1$). اِحسب \omterm{def:g8:speed:average}{المتوسط
المرجّح}. وأيّ متوسط بسيط (غير مرجّح) كانت ستحصل عليه، ولماذا
يختلفان؟
\end{exercise}

\begin{exercise}[$\star\star$]\label{exo:g8:speed:10}
تسير سيارة $120$ كم \omterm{def:g8:speed:def}{بسرعة} $80$ كم/سا ثم $120$ كم \omterm{def:g8:speed:def}{بسرعة} $120$ كم/سا.
\begin{enumerate}
  \item اِحسب مدّة كل مرحلة، ثم \omterm{def:g8:speed:def}{السرعة المتوسطة} على
        الرحلة كلها.
  \item واِشرح لماذا يكون الجواب أصغر من $100$ كم/سا، وهي منتصف
        السرعتين.
\end{enumerate}
\end{exercise}

\begin{exercise}[$\star\star\star$]\label{exo:g8:speed:11}
تبعد قريتان $18$ كم إحداهما عن الأخرى. تغادر آية الأولى عند $10{:}00$
ماشيةً \omterm{def:g8:speed:def}{بسرعة} $5$ كم/سا نحو الثانية؛ ويغادر بدر الثانية في
الوقت نفسه ماشيًا \omterm{def:g8:speed:def}{بسرعة} $4$ كم/سا نحو الأولى. ففي أيّ ساعة
يلتقيان، وعلى أيّ بُعد من قرية آية؟ (فهما معًا
يقلّصان الفجوة \omterm{def:g8:speed:def}{بسرعة} $5 + 4$ كم/سا.)
\end{exercise}

\section{مسألة: المتوسط التوافقي، أو لماذا تفسد رحلة العودة
المتوسط}

\begin{problem}[{مسألة نهاية الأسبوع --- \omterm{def:g8:speed:def}{السرعة المتوسطة} على مسافتين
متساويتين هي $\frac{2 v_1 v_2}{v_1 + v_2}$، وهي لا تتجاوز أبدًا
منتصف السرعتين}]
\label{pb:g8:speed:1}
اِنتهى كل من \cref{ex:g8:speed:twolegs} و \cref{exo:g8:speed:10}
بالمفاجأة نفسها: اِذهب \omterm{def:g8:speed:def}{بسرعة}، وعُد \omterm{def:g8:speed:def}{بسرعة}
أخرى، فتقع \omterm{def:g8:speed:def}{السرعة المتوسطة} \emph{تحت} منتصف
السرعتين. وتجد هذه المسألة الصيغة المضبوطة وراء المفاجأة
--- وهي نوع جديد من المتوسطات، هو \emph{المتوسط
التوافقي}\index{متوسط توافقي} --- وتبرهن على أن المفاجأة
مبرهنة، وتدفعها إلى حدّها المنطقي: أي لحاق لا تستطيعه
أيّ \omterm{def:g8:speed:def}{سرعة} في العالم.

\medskip
\noindent\textbf{الجزء الأول --- صيغة رحلة الذهاب والعودة.}
رحلة تقطع مسافة $d$ \omterm{def:g8:speed:def}{بسرعة} $v_1$، ثم المسافة
$d$ نفسها \omterm{def:g8:speed:def}{بسرعة} $v_2$ (وكلها أعداد موجبة).
\begin{enumerate}
  \item عبّر عن المدّتين، ثم عن المدّة الكلية، بكسور
        تتضمّن $d$ و $v_1$ و $v_2$.
  \item \omterm{def:g8:speed:def}{السرعة المتوسطة} للرحلة كلها هي
        $\bar v = \dfrac{2d}{\ \dfrac{d}{v_1} + \dfrac{d}{v_2}\ }$
        --- وهو \omterm{def:g6:fractions:def}{كسر} ذو طابقين
        (\cref{ex:g8:fractions:complex}). بسّطه وبَرهِن على أن
        \[
        \bar v = \frac{2\, v_1 v_2}{v_1 + v_2} .
        \]
  \item تحقّق من الصيغة بمقابلتها بالحسابين المذكورين أعلاه:
        $v_1 = 30$ و $v_2 = 15$ كم/سا
        (\cref{ex:g8:speed:twolegs})، و $v_1 = 80$ و
        $v_2 = 120$ كم/سا (\cref{exo:g8:speed:10}).
  \item وقد اِختفت المسافة $d$ من الصيغة. فماذا يعني ذلك،
        عمليًّا، بالنسبة إلى الدراج؟
  \item والآن اِفترض أن الرحلة تقضي بدل ذلك الزمن
        $T$ نفسه بكل \omterm{def:g8:speed:def}{سرعة}. بيّن أن \omterm{def:g8:speed:def}{السرعة
        المتوسطة} تصير حينئذٍ $\frac{v_1 + v_2}{2}$، أي المنتصف البسيط.
        وبلغة \cref{def:g8:speed:average}: بأيّ
        مقدار يجب أن يُرجَّح متوسط السرعات دائمًا،
        وما الذي يتغيّر بين السيناريوهين؟
\end{enumerate}

\medskip
\noindent\textbf{الجزء الثاني --- التوافقي مقابل الحسابي.}
من أجل عددين موجبين $v_1$ و $v_2$، ضع
\[
H = \frac{2\, v_1 v_2}{v_1 + v_2}
\quad\text{(وهو \emph{متوسطهما التوافقي})،}
\qquad
A = \frac{v_1 + v_2}{2}
\quad\text{(وهو \emph{متوسطهما الحسابي}).}
\]
\begin{enumerate}[resume]
  \item اِحسب $H$ و $A$ من أجل الزوجين $(30, 15)$ و
        $(80, 120)$. فأيّ المتوسطين يفوز في كل مرة؟
  \item ضع $A - H$ على \omterm{def:g4:fractions:def}{المقام} المشترك $2(v_1 + v_2)$،
        واُنشر \omterm{def:g4:fractions:def}{البسط} بالمتطابقات الشهيرة
        (\cref{pb:g8:equations:1})، وبَرهِن على أن
        \[
        A - H = \frac{(v_1 - v_2)^2}{2\,(v_1 + v_2)} .
        \]
  \item اِستنتج المبرهنة الكامنة وراء كل المفاجآت: من أجل السرعات
        الموجبة، $H \leq A$ دائمًا، مع التساوي بالضبط حين
        $v_1 = v_2$. ولماذا يحسم مربع في \omterm{def:g4:fractions:def}{البسط}
        الأمر؟
  \item بَرهِن على المتطابقة الأنيقة $H \times A = v_1 \times v_2$:
        فالمتوسطان يعودان بالضرب إلى \omterm{def:g2:mult:def}{الجداء} الأصلي.
  \item تقطع سيارة $60$ كم \omterm{def:g8:speed:def}{بسرعة} $40$ كم/سا و $60$ كم \omterm{def:g8:speed:def}{بسرعة}
        $60$ كم/سا. توقّع \omterm{def:g8:speed:def}{السرعة المتوسطة} بالصيغة،
        ثم أكّدها بالطريق الطويل، بالأزمنة والمسافة
        الكلية.
\end{enumerate}

\medskip
\noindent\textbf{الجزء الثالث --- اللحاق المستحيل.}
\begin{enumerate}[resume]
  \item تدخل دراجة سباقًا طوله $30$ كم آملةً أن يكون متوسطها
        $30$ كم/سا. فكم من الوقت يمكن أن يستغرق السباق على
        الأكثر؟ وهي تقطع أول $15$ كم \omterm{def:g8:speed:def}{بسرعة} $15$ كم/سا: فكم
        يبقى من ميزانية الوقت ذاك؟
  \item بيّن أن المتوسط المرجوّ صار مستحيلًا
        حرفيًّا: اِحسب سرعتها المتوسطة إن قطعت
        النصف الثاني \omterm{def:g8:speed:def}{بسرعة} $90$ كم/سا، واِشرح لماذا تبقى دون
        $30$ كم/سا حتى \omterm{def:g8:speed:def}{بسرعة} هائلة كيفما كانت.
  \item ثم تخفض هدفها إلى $20$ كم/سا. فأيّ \omterm{def:g8:speed:def}{سرعة} على
        مسافة $15$ كم الثانية تحقّقه؟ تحقّق من جوابك بصيغة
        \omterm{pb:g8:speed:1}{المتوسط التوافقي}.
  \item حنفية باردة وحدها تملأ حوضًا في $20$ دق؛ وحنفية ساخنة وحدها
        تملؤه في $30$ دق. فأيّ \omterm{def:g6:fractions:def}{كسر} من الحوض تملؤه الحنفيتان
        معًا في الدقيقة، وكم من الوقت تستغرقان؟
        وقارن جوابك بالمقدار $H(20, 30)$ --- فماذا
        تلاحظ؟
  \item تطير طائرة المسار نفسه ذهابًا وإيابًا: $900$ كم/سا
        مع الريح، و $700$ كم/سا ضدّها. اِحسب
        \omterm{def:g8:speed:def}{السرعة المتوسطة} للرحلة كلها، واِشرح في جملة واحدة لماذا
        تُطيل الريح، مهما هبّت، رحلة الذهاب
        والعودة دائمًا.
\end{enumerate}
\end{problem}
